\documentclass[11pt,a4paper]{article}

\usepackage[utf8]{inputenc}
\usepackage[T1]{fontenc}
\usepackage[spanish]{babel}
\usepackage[margin=2cm, bottom=6.2cm, footskip=1.7cm]{geometry}
\usepackage{graphicx}
\usepackage{array}
\graphicspath{{../../01_Propuesta/Logos/}{./}}
\usepackage{fancyhdr}
\usepackage{xcolor}
\usepackage[hidelinks]{hyperref}
\usepackage{enumitem}
\setlist{nosep}
\setlength{\parindent}{0pt}
\setlength{\parskip}{4pt}

\definecolor{verdeosc}{RGB}{26,107,60}
\definecolor{verdemid}{RGB}{46,139,87}

\pagestyle{fancy}
\fancyhf{}
\fancyfoot[C]{%
  \noindent{\color{verdeosc}\rule{\textwidth}{0.8pt}}\\[6pt]
  \tiny
  \hyphenpenalty=10000\exhyphenpenalty=10000
  \noindent
  \begin{minipage}{\textwidth}
  \begin{minipage}[t]{4.4cm}
  {\color{verdeosc}\textbf{ORGANIZA}}\\[3pt]
  Grupo de Investigación IN SILICO\\[6pt]
  {\color{verdeosc}\textbf{CO-ORGANIZADORES}}\\[3pt]
  Departamento de Biología y Química\\
  Departamento de Ingeniería Agroindustrial\\
  Departamento de Ingeniería Agrícola\\
  Departamento de Física\\
  Universidad de Sucre
  \end{minipage}%
  \hfill
  \begin{minipage}[t]{3.7cm}
  {\color{verdeosc}\textbf{COORDINACIÓN}}\\[3pt]
  \textbf{Aldo F. Combariza}\\
  Coordinación General\\
  \textbf{María Ximena Díaz}\\
  Coordinación Académica\\
  \textbf{Sebastián Vargas}\\
  Coordinación Logística\\
  \textbf{Selena Arias Avila}\\
  Coordinación Administrativa
  \end{minipage}%
  \hfill
  \begin{minipage}[t]{3.7cm}
  {\color{verdeosc}\textbf{CONTACTO}}\\[3pt]
  \textbf{SICVEC 2026}\\
  19--20 de octubre de 2026\\
  C.C. Guacarí, Sincelejo\\
  Modalidad híbrida\\
  \href{mailto:insilico@unisucre.edu.co}{insilico@unisucre.edu.co}
  \end{minipage}%
  \end{minipage}%
}
\renewcommand{\headrulewidth}{0pt}

\begin{document}

\noindent
\includegraphics[height=1.6cm]{logo_sicvec2026.png}%
\hspace{10pt}{\color{verdeosc}\vrule width 1.2pt height 1.5cm}\hspace{10pt}%
\includegraphics[height=1.28cm]{Logo_insilico.jpeg}%
\hspace{10pt}{\color{verdeosc}\vrule width 1.2pt height 1.5cm}\hspace{10pt}%
\includegraphics[height=1.3cm]{Logo_Bio.png}\hspace{6pt}%
\includegraphics[height=1.3cm]{logo_agroin.png}\hspace{6pt}%
\includegraphics[height=1.3cm]{logo_ingagrocola.png}\hspace{6pt}%
\includegraphics[height=1.3cm]{logo_fisica_lifi.png}%
\hspace{10pt}{\color{verdeosc}\vrule width 1.2pt height 1.5cm}\hspace{10pt}%
\includegraphics[height=1.3cm]{logo_unisucre.png}

\vspace{0.15cm}

\noindent
\textbf{Sincelejo, 8 de octubre de 2026}

\vspace{0.15cm}

\noindent
Señor\\
Guillermo Alejandro Corrales Ducuara\\
Estudiante, Maestría en Tecnologías Limpias\\
Universidad Santo Tomás, sede Bogotá D.C.\\
Bogotá D.C., Colombia

\vspace{0.1cm}

\noindent
\textbf{\large Asunto: Confirmación de su participación como ponente — SICVEC 2026}

\vspace{0.1cm}

Estimado Guillermo:

\vspace{0.1cm}

Muchas gracias por su mensaje y por aceptar la invitación a participar como \textbf{ponente} en el \textbf{Simposio Internacional de Ciencia Verde y Economía Circular (SICVEC 2026)}. Será un gusto contar con su trabajo sobre metano en codigestión anaeróbica, dentro de su formación en Tecnologías Limpias en la Universidad Santo Tomás.

Para incluir su charla en el programa, le pedimos enviarnos lo siguiente:

\noindent
\textbf{INFORMACIÓN REQUERIDA}

\begin{enumerate}
  \item Título definitivo de la ponencia
  \item Un breve resumen (sin límite de palabras)
  \item Una reseña corta suya (biografía)
  \item Una foto suya para la promoción del evento
\end{enumerate}

No hay un formato obligatorio; puede enviarlo en el cuerpo del correo o adjunto en el formato que le resulte más cómodo.

\vspace{0.1cm}

En los próximos días le compartiremos el cronograma definitivo con la franja horaria asignada para su ponencia, junto con las instrucciones técnicas para el evento.

\vspace{0.1cm}

Quedamos atentos a su información. Nuevamente, muchas gracias por su disposición.

\vspace{0.15cm}

Atentamente,

\vspace{0.15cm}

\textbf{Dr. Aldo F. Combariza}\\
Coordinador General, SICVEC 2026\\
Grupo de Investigación IN SILICO · Departamento de Biología y Química\\
Facultad de Educación y Ciencias · Universidad de Sucre\\
\href{mailto:aldo.combariza@unisucre.edu.co}{aldo.combariza@unisucre.edu.co}

\vspace{0.2cm}

\end{document}
